\section{Introduction}

A humanoid that climbs onto a waist-high box uses its hands, knees, and shins, in an order that is hard to find by trial and error and easy to see in a person doing it. Two families of methods turn that observation into a controller, and each gives up something. \emph{Trackers} follow a retargeted clip frame by frame~\cite{deepmimic,beyondmimic,zest,omniretarget}. They reproduce the motion, but they need the clip at deployment and fail when the robot does not start where the clip does. \emph{Goal-driven} policies are told only where to end up, so they adapt to where they start, but a reward that says only ``you stand on the box'' gives no signal until the robot stands on the box, and the shaped rewards that replace it are engineered per behaviour.

Hybrid methods train both objectives at once, in parallel worlds of one simulation that share a policy~\cite{hil,wang2026generalizing}. Their results agree on two points. The mixture matters enormously: with a sparse goal reward, removing the reference worlds takes success on box jumps and climbs from 1.00 to 0.00~\cite{wang2026generalizing}, and adding goal worlds before the policy can follow the references does the same. And the mixture is set by hand: a fixed number of samples of tracking followed by an even split~\cite{hil}, or a schedule driven by how well the policy tracks~\cite{wang2026generalizing}. Neither measures what the reference worlds do for the goal task, which is the reason they are there.

This paper asks what the worlds of each kind \emph{cause}, measures it inside the training simulation, and allocates worlds by it. We call the principle \emph{credit the cause}: a kind of world receives worlds for the failures of the evaluation that it has been measured to remove, in the tasks where it is able to remove them. The system has three kinds of worlds (Fig.~\ref{fig:system}). \emph{Reference worlds} start on a frame of a clip and imitate a tracker. \emph{Goal worlds} start on a frame of a clip or away from it and receive one Boolean at the end of the episode: whether the robot stands at the goal. \emph{Evaluation worlds} run the policy without exploration noise on every task from a fixed ladder of starts and give nothing to the learner. Their Booleans are the only signal the allocation reads.

We study the Unitree G1 climbing onto and over boxes of 0.38 to 1.14\,m, with the clips and boxes of a public dataset~\cite{omniretarget}. The policy we train never sees a clip: it observes the robot, a height scan of the terrain, and the pose of the goal. Our contributions are:
\begin{itemize}
\item A measurement of what each kind of world causes. Goal worlds alone learn nothing in the budget we tested, not even to stand. A tracker reaches the goal from the first frame of its clip and almost never from starts half a metre away. Reference worlds alone give a clip-free policy the motion but not the adaptation. Goal worlds added to them give both (Sec.~\ref{sec:cause}).
\item Evidence that the share of worlds of each kind is a strong and one-sided lever: with a tenth of the worlds on references nothing is learned, with all of them the policy stays tied to the clip, and the best share moves during training (Sec.~\ref{sec:share}).
\item An allocation rule that reads only the evaluation worlds and costs nothing. It needs no schedule, and reaches our success levels in 1.7 times fewer samples than the best fixed share found by a sweep and 2.0 to 2.4 times fewer than two published hand schedules (Sec.~\ref{sec:schedules}). Two learners that try neighbouring shares in parallel measure the same direction directly, at twice the samples.
\item One policy for \tasksall{} box-climbing tasks from public data. It solves 103 of them from the starts of the clips and 104 from displaced starts, against 62 and one for the tracker that teaches it, and 41 tasks that the tracker does not solve at all (Sec.~\ref{sec:all}).
\end{itemize}
All results are in simulation; we state what we did not test in Sec.~\ref{sec:limits}.

\section{Related Work}

\textbf{Tracking and its data.} DeepMimic~\cite{deepmimic} established reward-driven tracking of a reference with starts on its frames. Recent trackers scale it to hardware and to contact-rich skills~\cite{beyondmimic,zest}. OmniRetarget~\cite{omniretarget} generates clips of a G1 on terrain that preserve the contacts of the human, and releases those we use. A tracker needs its reference to act. Distillation removes it at deployment by training a student on the tracker's actions under a sparser interface~\cite{maskedmimic,hover}, the DAgger scheme~\cite{dagger}; the student is then trained only on what the tracker visits.

\textbf{Imitation with goals.} Adversarial priors add a style reward to a task reward~\cite{amp,amphardware}; on a parkour course, starting them from scratch completes 0.11 of the courses against 0.85 after a tracking stage~\cite{hil}. HIL~\cite{hil} runs tracking worlds and goal worlds in parallel with one observation space, after four billion samples of tracking alone. Wang et al.~\cite{wang2026generalizing} train a G1 on box jumps and climbs with imitation worlds and goal worlds whose reward is sparse, one policy per behaviour; a scalar that follows tracking success lowers the share of imitation worlds from 1.0 to 0.5 and removes an assistive force. ConsMimic~\cite{consmimic} adapts the weight of a style reward to keep the task return near its best, and notes that it cannot tell which parts of the demonstrations help. We keep the parallel worlds of these methods, supervise the reference worlds with a tracker instead of a tracking reward alone, and set the split from what evaluation worlds measure.

\textbf{Curricula and samplers.} Success Guided Sampling~\cite{sgs} samples a fixed set of task configurations by their own success rate, with a Boolean terminal reward and resets that reach near the goal~\cite{omnireset}; its ablations show little sensitivity to the parameters of the sampling rule and no learning without the resets near the goal. Learning-progress curricula score a task by the change of its own return~\cite{lpacrl,graves2017automated}. Starts on demonstration states turn a sparse reward into a curriculum~\cite{florensa2017reverse,demostart,scout,goexplore}. These methods choose among tasks or starts of one objective by statistics of the training episodes. We allocate between kinds of supervision, and read held-out evaluation episodes instead. Population-based training~\cite{pbt} selects among learners by measured performance; our two trial learners are its smallest instance, applied to the share of worlds.

\textbf{Humanoid and legged parkour.} Perceptive parkour policies are trained with shaped rewards and terrain curricula~\cite{zhuang2023robot,cheng2024extreme,hoeller2024anymal,zhuang2024humanoid}. Our goal reward is one Boolean, and what the shaped terms would provide comes from the clips.

